%\section{Introducción}
%El monitoreo automatizado de la fauna silvestre mediante imágenes aéreas es crucial para la conservación ecológica, pero enfrenta un desafío crítico en la sabana africana: la detección precisa en escenarios de alta densidad. Especies gregarias como elefantes y búfalos forman manadas compactas con severa oclusión y superposición entre individuos, condiciones donde los sistemas de visión por computador convencionales fallan sistemáticamente.
%
%El problema fundamental radica en las arquitecturas tradicionales basadas en CNN y anclas, que dependen del algoritmo de Supresión de No Máximos (NMS). En manadas densas, el NMS elimina erróneamente detecciones válidas de animales contiguos, provocando un subconteo crónico. Aunque alternativas como HerdNet \cite{herdnet} evitan el NMS mediante la estimación de densidad por puntos, su dependencia de backbones CNN limita la captura del contexto global necesario para una clasificación robusta, afectando especialmente a las especies minoritarias.
%
%Este trabajo se justifica por la necesidad de una solución que supere simultáneamente las limitaciones de oclusión y discriminación semántica. El objetivo principal es implementar y evaluar la arquitectura \textbf{RF-DETR} (\textit{Real-Time Detection Transformer}) \cite{rf-detr} en este contexto desafiante. Nuestra hipótesis es que el mecanismo de atención de los Transformers, capaz de modelar dependencias globales, junto con la predicción de conjuntos \textit{end-to-end} que elimina por completo la necesidad del NMS, ofrece una solución estructural al problema del subconteo en manadas.
%
%Las principales contribuciones de este estudio son:  (1) Una comparación de tres variantes de RF-DETR (Nano, Small, Large) frente al baseline HerdNet. (2) Un análisis del compromiso entre precisión en alta densidad y costo computacional de las arquitecturas, identificando configuraciones óptimas para el despliegue práctico. (3) La validación de una estrategia de entrenamiento de dos fases con Minería de Negativos Difíciles (HNM) para refinar la precisión en fondos complejos.

\section{Introduction}
Automated wildlife monitoring using aerial imagery is crucial for ecological conservation, but faces a critical challenge in the African savanna: precise detection in high-density scenarios. Gregarious species such as elephants and buffalo form compact herds with severe occlusion and overlap between individuals, conditions where conventional computer vision systems systematically fail.

The fundamental problem lies in traditional CNN-based and anchor-based architectures, which rely on the Non-Maximum Suppression (NMS) algorithm. In dense herds, NMS erroneously eliminates valid detections of contiguous animals, causing chronic undercounting. Although alternatives such as HerdNet \cite{herdnet} avoid NMS through point-based density estimation, their dependence on CNN backbones limits the capture of global context necessary for robust classification, particularly affecting minority species.

This work is justified by the need for a solution that simultaneously overcomes the limitations of occlusion and semantic discrimination. The main objective is to implement and evaluate the \textbf{RF-DETR} (\textit{Real-Time Detection Transformer}) \cite{rf-detr} architecture in this challenging context. Our hypothesis is that the Transformer's attention mechanism, capable of modeling global dependencies, together with \textit{end-to-end} set prediction that completely eliminates the need for NMS, offers a structural solution to the undercounting problem in herds.

The main contributions of this study are: (1) A comparison of three RF-DETR variants (Nano, Small, Large) against the HerdNet baseline. (2) An analysis of the trade-off between precision in high density and computational cost of the architectures, identifying optimal configurations for practical deployment. (3) The validation of a two-phase training strategy with Hard Negative Mining (HNM) to refine precision in complex backgrounds.
